\documentclass[a4paper,12pt]{article}

\usepackage{titlesec}
\usepackage{hyperref}
\usepackage{amsmath}
\usepackage{amssymb}
\usepackage{braket}
\usepackage[utf8]{inputenc}
\usepackage{tikz}
\usepackage{cancel}
\usepackage[top=3cm, bottom=3cm, left=2cm, right=2cm]{geometry}
\usepackage{graphicx}
\usepackage{subcaption}
\usepackage{float}
\usetikzlibrary{arrows.meta, angles, quotes}

\allowdisplaybreaks

\titleformat{\section}
{\normalfont\normalsize\bfseries}
{\thesection}{1em}{}
\titleformat{\subsection}
{\normalfont\small\bfseries}
{\thesubsection}{1em}{}
\titleformat{\subsubsection}
{\normalfont\small\bfseries}
{\thesubsubsection}{1em}{}

\title{
	\vspace{2cm}
	\Huge \textbf{Titolo ufficiale} \newline
	\large frontespizio fatto bene
}

\author{\Large Francesco Ciavarella - University of Padua \\[10pt]}
 
\begin{document}
	
	\maketitle
	\thispagestyle{empty}
	\newpage
	
	\tableofcontents
	\newpage
	
	\section{Introduction}
	\label{sec:intro}
	Introduzione generale e filosofica sul lavoro di tesi; temi da toccare : importanza delle tecnologie quantistiche, formalismo della matrice densità, algoritmi per quantum computer, ...
	deve dare l'impressione che serva a qualcosa di pratico
	
	\clearpage
	\newpage
	
	\section{Open Quantum System and Quantum Master Equation}
	\label{sec:Open_Quantum_System_and_Quantum_Master_Equation}
	
	subsection 1 : Formalismo matrice densità vs Schrodinger (parlare di alcune metriche importanti tipo norma $l_1$, entropia VN, ... ); evoluzione tramite QME, mappa dinamica eccetera 
	subsection 2 : forma di riferimento Lindbald -> approfondisci.
	subsection 3 : unraveling di QME : KRAUS OPERATORS; SSE breve trattazione e magari accenno all'algoritmo QJ; Collisonal Methods -> trattati meglio nella sezione 3 
	
	\section{Collisional Methods} 
	\label{sec:Collisional_Methods}
	Trattazione matematica generica, focus su interpretazione fisica, es ancilla come effetto del'ambiente, modellabile per rappresentare env divesi, metto qui i due limiti diffusivo e quantum jump; 	CAPIRE COME SONO LEGATI A MISURE DIVERSE !! 
	subsection 1 :
	sarebbe bello parlare della differenza intrinseca tra traiettorie originate con schemi di unraveling diversi, itroducendo alcune metriche utilizzate per carterizzare queste differenze : ad esempio la probabilità itrodotta nei kraus operators, la fidelity tra step precedente e successivo, quindi tra stati puri, come valore dell'intensità di cambiamento, ... 
	(VERIFICA PROBABILITÀ BASSA PER EVENTI RARI COME QJ)
	subsection 1 :
	trattare i temi di convergenza finite sampling,
	A parte le differenze intrinseche ci sono poi le importanti differnze di convergenza, quindi si deve parlare di TLC e varianza !!
	\textbf{qui dividere in statoo dell'arte e obiettivo della mia tesi}.
	\section{First Case of Study : the Exciton Dimer}
	\label{sec:Exciton_Dimer}
	Breve introduzione sul tipo di sistema, le caratteristiche, la dinamica di pure dephasing studiata (associata a ENAQT), la lindblad di riferimento con anche plot e perchè tende allo stato puramente misto.
	\subsection{Quantum Jump Limit} 
	\label{sub:QJ_limit_Exc}
	trattazione completa con formule e derivazione, prese in modo selettivo da MC; risultati di convergenza, singole traiettorie, ...
	(metto analisi della convergenza ache per dt diversi una volta per tutte in modo da decidere che userò dt = 0.1 sufficientemente piccolo)
	\subsection{Diffusive Limit} 
	\label{sub:Diffusive_limit_Exc}
	trattazione completa con formule e derivazione, prese in modo selettivo da MC; risultati di convergenza, singole traiettorie, ...
	
	\subsection{Intermediate Regime} 
	\label{sub:Intermediate_Regime_Exc}
	trattazione completa con formule e derivazione, prese in modo selettivo da MC; risultati di convergenza, singole traiettorie, ...
	
	\subsection{Convergence Analysis}
	\label{sub:Convergence_Analysis_Exc}
	Tutti i risultati sull'analisi della convergenza : Trace Distance e Fidelity come parametri di misura dell'errore, plot del fitting e verica che rispetta il TLC. 
	Analisi della varianza di un osservabile, in particolare i valro di aspettazioen delle componenti del vetore di bloch (caratterizzano completamente il nostro sistema = ricostruzione matrice densità)
	
	\section{Second Case of Study : the Photoemission Qubit}
	\label{sec:Photoemission_Qubit}
	capire che cazzo fare della mia vita (soluzione torre di archimede)
	
	Breve intro sul tipod di problema da affronatre adesso, lindbald associata, ... ci starebbe un piccolo excursus sul problema affrontato da Coccia tramite SSE e del fatto che non salta mai
	
	Poi stessa trattazione già vista precedentemente
	\subsection{Quantum Jump Limit}
	\label{sub:QJ_limit_Photo}
	...
	
	\subsection{Diffusive Limit} 
	\label{sub:Diffusive_limit_Photo}
	...
	
	\subsection{Intermediate Regime} 
	\label{sub:Intermediate_Regime_Photo}
	...
	
	\subsection{Convergence Analysis}
	\label{sub:Convergence_Analysis_Photo}
	...
	
\end{document}